\section{まとめと今後の課題}

本ノートでは，格子ボルツマン法による変化球軌道シミュレーションの理論的背景（\secref{aerodynamics}--\secref{turbulence}）と実装方法（\secref{geometry-boundary}--\secref{gpu}），および検証結果（\secref{validation}）をまとめた。現時点で得られている主要な結論は次のとおりである。

\begin{itemize}
	\item 中心モーメント衝突演算子と\(\omega_{odd}\)/\(\omega_{higher}\)の分離により，本番レイノルズ数（\(\sim2\times10^5\)）で必要な緩和時間域でも数値的に安定な計算が可能であることを確認した。
	\item 周期境界における後流の巻き込みが，高レイノルズ数側のV\&V-2で見られた大きな誤差の主因であることを特定し，流入・流出境界の導入によってこれを解消した。
	\item 壁面条件（補間バウンスバック，実効壁位置）は複数の独立な検証により系統誤差が無いことを確認した一方，緩和時間に依存する6\%程度の誤差の原因は本稿執筆時点では未特定であり，コードの不確かさとして申告した上で本番実行に進んでいる。
	\item AA-pattern・局所グリッドリファインメント・GPU実装により，1条件あたりの計算時間を計画段階の見積り（15--20時間）から約1.3時間まで短縮した。
\end{itemize}

\subsection{今後の課題}

\begin{itemize}
	\item \textbf{緩和時間依存性の原因特定}：\secref{validation-highre}で残った6\%の緩和時間依存性について，曲面壁での補間バウンスバックと運動量交換法の組み合わせに特有の効果である可能性が高く，平板壁での検証（本稿の壁面条件の検証）だけでは十分条件になっていない。
	\item \textbf{グリッドリファインメントの多段化}：現状は1レベル・2倍細分化のみが実装されている。縫い目を1格子間隔以上で解像するには，細格子内へのさらなる入れ子（2レベル目）が必要である。
	\item \textbf{V\&V-3・V\&V-4の実施}：回転滑面球のマグヌス係数検証，および縫い目付き球体の固定・一様流中での検証を完了させる。
	\item \textbf{V\&V-5の実施}：スピン軸の向き・リリースポイントの確認またはキャリブレーション，「変化量」の定義の確定を経て，WBC2023決勝の検証ケースで本番実行と実測比較を行う。
	\item \textbf{可視化パイプラインの拡充}：解析モデル軌道の対話的可視化（段階1）は実装済みであり，CFD軌道の同様の可視化（段階2）が残っている。
\end{itemize}

本ノートは開発の進行に伴って随時更新する予定であり，日々の検証結果や設計判断のより詳細な記録は，本リポジトリの設計ノート（\texttt{docs/design/DESIGN.md}）を参照されたい。
